%! Vector-Valued Functions — Unit 4, Lesson 4.9 — Warm-Up
\documentclass[10pt]{article}
\usepackage{precalculus-article}
\usepackage{precalculus-boxes}
\newcommand{\TallMath}[1]{$\displaystyle #1\rule[-1.4em]{0pt}{3.2em}$}

\begin{document}

\pageheader{Unit 4, Lesson 4.9}{Warm-Up}
\namedateperiod

\vspace{0.12in}

\begin{notesbox}{Spiral Review}
\small

\textbf{1.}\quad Given the parametric function $f(t)=(2t,\;t^{2}-1)$, find $f(3)$.

\vspace{0.08in}
\hspace{1em} $x(3) = $ \blank{3cm} \qquad $y(3) = $ \blank{3cm}
\qquad $f(3) = $ \blank{3cm}

\vspace{0.22in}

\textbf{2.}\quad Find the magnitude of the vector $\langle 3,\,-4\rangle$.

\vspace{0.08in}
\hspace{1em} $\|\langle 3,\,-4\rangle\| = \sqrt{\blank{1.5cm}^{2}+(\blank{1.5cm})^{2}}
= \sqrt{\blank{2cm}} = $ \blank{2cm}

\vspace{0.22in}

\textbf{3.}\quad A particle is at position $(0,\,5)$ when $t=0$ and at position
$(6,\,5)$ when $t=3$.
Find the average rate of change of $x(t)$ over the interval $0\le t\le 3$.

\vspace{0.08in}
\hspace{1em} $\dfrac{\Delta x}{\Delta t} = \dfrac{\blank{1.8cm}-\blank{1.8cm}}{\blank{1.8cm}-\blank{1.8cm}}
= \dfrac{\blank{1.5cm}}{\blank{1.5cm}} = $ \blank{2cm}

\end{notesbox}

\end{document}
